\noindent\textbf{Exercise 1.1.14 (Dyadic approximation and Jordan measurability).}\par
Let $E\subset\mathbb{R}^d$ be bounded. For $n\geq 0$, let
$\mathcal D_n$ be the dyadic cubes of side length $h_n=2^{-n}$,
and set
\[
N_n^-:=\#\{Q\in\mathcal D_n:Q\subseteq E\},\qquad
N_n^+:=\#\{Q\in\mathcal D_n:Q\cap E\neq\varnothing\}.
\]
Show that $E$ is Jordan measurable if and only if
\[
2^{-dn}(N_n^+-N_n^-)\longrightarrow0.
\]
In this case, show that both $2^{-dn}N_n^-$ and $2^{-dn}N_n^+$
converge to $m(E)$.

\begin{proof}
For any bounded $T\subset\mathbb R^d$, write
\[
I_n(T):=\bigcup_{\substack{Q\in\mathcal D_n\\Q\subseteq T}}Q,
\qquad
O_n(T):=\bigcup_{\substack{Q\in\mathcal D_n\\Q\cap T\neq\varnothing}}Q.
\]
Only finitely many cubes at scale $n$ intersect $T$; hence
$I_n(T)$ and $O_n(T)$ are elementary. The cubes at each scale
partition $\mathbb R^d$, so
\[
I_n(T)\subseteq T\subseteq O_n(T),\qquad
m(I_n(E))=2^{-dn}N_n^-,\quad m(O_n(E))=2^{-dn}N_n^+.
\tag{1}
\]

\emph{Claim.} If $S$ is elementary, then
\[
m(I_n(S))\longrightarrow m(S),\qquad
m(O_n(S))\longrightarrow m(S).
\tag{2}
\]
Indeed, for an interval $J\subset\mathbb R$ of length $\ell$,
only the dyadic intervals meeting its two endpoints can contribute
to the difference between inner and outer approximations. Thus,
with $h=h_n$,
\[
(\ell-2h)_+\leq m_1(I_n(J))\leq\ell
\leq m_1(O_n(J))\leq\ell+2h.
\]
Consequently, both one-dimensional approximation measures tend
to $\ell$. For a box $R=J_1\times\cdots\times J_d$, the dyadic
partition is a product partition, so
\[
I_n(R)=\prod_{j=1}^d I_n(J_j),\qquad
O_n(R)=\prod_{j=1}^d O_n(J_j).
\]
The elementary product formula now gives
$m(I_n(R))\to |R|$ and $m(O_n(R))\to |R|$.
For an elementary set $S=\bigsqcup_{k=1}^M R_k$, note that
\[
\sum_{k=1}^M m(I_n(R_k))\leq m(I_n(S))\leq m(S)
\leq m(O_n(S))\leq\sum_{k=1}^M m(O_n(R_k)).
\]
Both outermost expressions converge to
$\sum_{k=1}^M|R_k|=m(S)$, proving (2) by squeezing.

Suppose first that $E$ is Jordan measurable. Given
$\varepsilon>0$, Exercise~1.1.5 supplies elementary sets
$A\subseteq E\subseteq B$ such that
$m(B\setminus A)<\varepsilon/3$.
Since dyadic inner and outer approximations preserve inclusions,
\[
I_n(A)\subseteq I_n(E)\subseteq O_n(E)\subseteq O_n(B).
\]
By (2), for all sufficiently large $n$,
\[
0\leq m(A)-m(I_n(A))<\varepsilon/3,
\]
and
\[
0\leq m(O_n(B))-m(B)<\varepsilon/3.
\]
Therefore
\[
\begin{aligned}
0\leq m(O_n(E))-m(I_n(E))
&\leq m(O_n(B))-m(I_n(A))\\
&=[m(O_n(B))-m(B)]+[m(B)-m(A)]\\
&\qquad+[m(A)-m(I_n(A))] < \varepsilon.
\end{aligned}
\]
By (1), $2^{-dn}(N_n^+-N_n^-)\to0$.

Conversely, suppose $2^{-dn}(N_n^+-N_n^-)\to0$.
For every $\varepsilon>0$ we may choose $n$ with
\[
m(O_n(E)\setminus I_n(E))
=m(O_n(E))-m(I_n(E))<\varepsilon.
\]
Since $I_n(E)\subseteq E\subseteq O_n(E)$ and both approximations
are elementary, Exercise~1.1.5 implies that $E$ is Jordan measurable.

Finally, for a Jordan measurable $E$, monotonicity gives
\[
m(I_n(E))\leq m(E)\leq m(O_n(E)).
\]
The difference of the outer bounds tends to zero, so both
converge to $m(E)$. Apply (1) to obtain
\[
\displaystyle
m(E)=\lim_{n\to\infty}2^{-dn}N_n^-
=\lim_{n\to\infty}2^{-dn}N_n^+.
\]
\end{proof}
